\documentclass[UTF8]{article}

% ========== 基础宏包 ==========
\usepackage{ctex}                % 中文支持
\usepackage{amsmath,amssymb}     % 数学公式
\usepackage{geometry}            % 页面边距
\usepackage{titlesec}            % 自定义标题格式
\usepackage{enumitem}            % 列表格式控制
\usepackage{microtype}           % 微排版优化
\usepackage{xcolor}              % 颜色支持
\usepackage{tcolorbox}           % 彩色盒子（核心美化包）
\tcbuselibrary{most}             % 加载大多数 tcolorbox 库

% ========== 页面设置 ==========
\geometry{a4paper, margin=2.5cm}
\linespread{1.2}                 % 1.2倍行距
\setlength{\parskip}{0.5ex}      % 段落间距

% ========== 标题格式 ==========
\titleformat{\section}{\Large\bfseries\centering}{\thesection}{1em}{}
\titleformat{\subsection}{\normalsize\bfseries}{\thesubsection}{0.8em}{}
\titleformat{\subsubsection}{\normalsize\itshape}{\thesubsubsection}{0.8em}{}

% ========== 列表环境 ==========
\setlist[itemize]{leftmargin=*, nosep, topsep=0.3ex, parsep=0ex, partopsep=0ex}
\setlist[enumerate]{leftmargin=*, nosep, topsep=0.3ex, parsep=0ex, partopsep=0ex}

% ========== 自定义彩色模块（tcolorbox） ==========
% ========== 自定义彩色模块（tcolorbox） ==========
\newtcolorbox{definitionbox}{
    colback=blue!5!white,
    colframe=blue!60!black,
    fonttitle=\bfseries,
    title=定义,
    boxrule=0.8pt,
    arc=4pt,
    left=8pt,
    right=8pt,
    top=6pt,
    bottom=6pt,
    drop fuzzy shadow=blue!10!white,
    before skip=6pt plus 2pt,
    after skip=6pt plus 2pt
}

\newtcolorbox{theorembox}{
    colback=orange!5!white,
    colframe=orange!70!black,
    fonttitle=\bfseries,
    title=核心定理,
    boxrule=0.8pt,
    arc=4pt,
    left=8pt,
    right=8pt,
    top=6pt,
    bottom=6pt,
    drop fuzzy shadow=orange!10!white,
    before skip=6pt plus 2pt,
    after skip=6pt plus 2pt
}

\newtcolorbox{proofbox}{
    breakable,
    enhanced,
    colback=green!5!white,
    colframe=green!50!black,
    fonttitle=\bfseries,
    title=证明,
    boxrule=0.8pt,
    arc=4pt,
    left=8pt,
    right=8pt,
    top=6pt,
    bottom=6pt,
    drop fuzzy shadow=green!10!white,
    before skip=6pt plus 2pt,
    after skip=6pt plus 2pt,
    pad at break*=2pt,
    skin first=enhanced,
    skin middle=enhanced,
    skin last=enhanced
}

\newtcolorbox{remarkbox}{
    colback=purple!5!white,
    colframe=purple!60!black,
    fonttitle=\bfseries,
    title=备注,
    boxrule=0.8pt,
    arc=4pt,
    left=8pt,
    right=8pt,
    top=6pt,
    bottom=6pt,
    drop fuzzy shadow=purple!10!white,
    before skip=6pt plus 2pt,
    after skip=6pt plus 2pt
}

% ========== 文档信息 ==========
\title{\textbf{强化学习课程笔记整理}}
\author{}
\date{}

\begin{document}

\maketitle
\thispagestyle{empty}

% ============================================================
\section{第一讲：强化学习基本概念}
% ============================================================

强化学习的核心框架基于马尔可夫决策过程，包含以下基础构成要素：

\begin{definitionbox}
  \textbf{状态（State）}：状态是智能体（Agent）对自身所处环境的表征。
  \begin{itemize}
    \item \textbf{状态空间}（State Space）：所有可能状态构成的集合，记为 $\mathcal{S}$；任意单个状态记为 $s$，满足 $s \in \mathcal{S}$。
  \end{itemize}
\end{definitionbox}

\begin{definitionbox}
  \textbf{动作（Action）}：智能体在特定状态下可执行的操作。
  \begin{itemize}
    \item 状态 $s$ 对应的动作空间记为 $\mathcal{A}(s)$；任意单个动作记为 $a$，满足 $a \in \mathcal{A}(s)$。
  \end{itemize}
\end{definitionbox}

\begin{definitionbox}
  \textbf{状态转移（State Transition）}：智能体在状态 $s$ 执行动作 $a$ 后，环境转移到下一状态的随机过程，由条件概率描述：
  \[
    p(s' \mid s, a) = \mathbb{P}(S_{t+1}=s' \mid S_t=s, A_t=a)
  \]
  其中 $S_t$、$A_t$ 分别表示第 $t$ 决策步的状态与动作。
\end{definitionbox}

\begin{definitionbox}
  \textbf{策略（Policy）}：策略定义了智能体的行为规则，即每个状态下的动作选择方式，记为 $\pi$。
  \begin{itemize}
    \item 随机策略表示为条件概率分布 $\pi(a \mid s)$，含义为：处于状态 $s$ 时选择动作 $a$ 的概率，满足归一性
    \[
      \sum_{a \in \mathcal{A}(s)} \pi(a \mid s) = 1 .
    \]
  \end{itemize}
\end{definitionbox}

\begin{definitionbox}
  \textbf{奖励（Reward）}：奖励是环境向智能体反馈的标量信号，是强化学习的优化目标来源。
  \begin{itemize}
    \item 即时奖励：在状态 $s$ 执行动作 $a$ 后获得的期望奖励记为 $r(s,a)$，即
    \[
      r(s,a) = \mathbb{E}\left[ R_{t+1} \mid S_t=s, A_t=a \right],
    \]
    其中 $R_{t+1}$ 是第 $t+1$ 步的随机奖励变量。
  \end{itemize}
\end{definitionbox}

\begin{definitionbox}
  \textbf{折扣回报（Discounted Return）}：回报是从当前时刻起，未来所有奖励的累积和；引入折扣因子以平衡即时奖励与长期奖励。
  \begin{itemize}
    \item 折扣因子 $\gamma \in [0,1)$，保证无限步交互下回报有限，取值越小表示智能体越看重即时奖励。
    \item 第 $t$ 步的折扣回报定义为：
    \[
      G_t = R_{t+1} + \gamma R_{t+2} + \gamma^2 R_{t+3} + \dots = \sum_{k=0}^{\infty} \gamma^k R_{t+k+1}.
    \]
    回报 $G_t$ 是依赖于轨迹的随机变量。
  \end{itemize}
\end{definitionbox}

\begin{definitionbox}
  \textbf{回合与任务分类}：
  \begin{itemize}
    \item \textbf{回合（Episode）}：具有明确终止状态的有限长度交互轨迹。
    \item \textbf{回合式任务（Episodic Tasks）}：任务存在终止状态，交互过程可划分为独立的回合。
    \item \textbf{持续性任务（Continuing Tasks）}：任务无终止状态，交互过程无限持续。
  \end{itemize}
\end{definitionbox}

\begin{definitionbox}
  \textbf{马尔可夫决策过程（MDP）}：马尔可夫决策过程是强化学习的标准数学模型，核心满足\textbf{马尔可夫性}：下一状态仅由当前状态和动作决定，与历史状态无关。
  \[
    \mathbb{P}(S_{t+1} \mid S_t, A_t, S_{t-1}, A_{t-1}, \dots) = \mathbb{P}(S_{t+1} \mid S_t, A_t).
  \]
  完整的 MDP 由五元组 $(\mathcal{S}, \mathcal{A}, p, r, \gamma)$ 构成：
  \begin{itemize}
    \item 状态空间 $\mathcal{S}$
    \item 动作空间 $\mathcal{A}$
    \item 状态转移概率函数 $p(s' \mid s,a)$
    \item 期望奖励函数 $r(s,a)$
    \item 折扣因子 $\gamma$
  \end{itemize}
\end{definitionbox}

\newpage

% ============================================================
\section{第二讲：贝尔曼方程}
% ============================================================

贝尔曼方程是强化学习的核心方程，刻画了价值函数的递归关系，是策略评估与优化的基础。

\subsection*{1. 状态价值函数}
在策略 $\pi$ 下，状态价值函数定义为：从状态 $s$ 出发，始终遵循策略 $\pi$ 交互所能获得的期望折扣回报：
\[
  v_\pi(s) = \mathbb{E}_\pi\left[ G_t \mid S_t = s \right].
\]
\begin{itemize}
  \item $v_\pi(s)$ 是状态 $s$ 的函数，取值完全依赖于策略 $\pi$；
  \item 价值函数本质是回报的期望，用于量化不同状态的长期优劣程度。
\end{itemize}

\subsection*{2. 贝尔曼期望方程推导}
利用回报的递归分解关系 $G_t = R_{t+1} + \gamma G_{t+1}$，代入状态价值函数的定义展开：
\begin{align*}
  v_\pi(s) &= \mathbb{E}_\pi\left[ G_t \mid S_t = s \right] \\
           &= \mathbb{E}_\pi\left[ R_{t+1} + \gamma G_{t+1} \mid S_t = s \right] \\
           &= \mathbb{E}_\pi\left[ R_{t+1} \mid S_t = s \right] + \gamma \cdot \mathbb{E}_\pi\left[ G_{t+1} \mid S_t = s \right].
\end{align*}

对第一项（即时奖励期望）按策略与状态转移展开：
\[
  \mathbb{E}_\pi\left[ R_{t+1} \mid S_t = s \right] = \sum_{a \in \mathcal{A}(s)} \pi(a \mid s) \cdot r(s,a).
\]

对第二项（未来回报期望），由马尔可夫性可得 $\mathbb{E}_\pi\left[ G_{t+1} \mid S_t=s, S_{t+1}=s' \right] = v_\pi(s')$，因此：
\[
  \mathbb{E}_\pi\left[ G_{t+1} \mid S_t = s \right] = \sum_{a \in \mathcal{A}(s)} \pi(a \mid s) \sum_{s' \in \mathcal{S}} p(s' \mid s,a) \cdot v_\pi(s').
\]

合并两项后得到\textbf{贝尔曼期望方程}：

\begin{theorembox}
  \textbf{贝尔曼期望方程（Bellman Expectation Equation）}
  \begin{equation} \label{eq:bellman_expect}
    v_\pi(s) = \sum_{a \in \mathcal{A}(s)} \pi(a \mid s) \left( r(s,a) + \gamma \sum_{s' \in \mathcal{S}} p(s' \mid s,a) \, v_\pi(s') \right), \quad \forall s \in \mathcal{S}.
  \end{equation}
  \textbf{直观含义}：当前状态的价值 = 即时奖励的期望 + 折扣后所有后续状态价值的加权期望。
\end{theorembox}

\subsection*{3. 矩阵向量形式}
对于有限离散状态空间，可将贝尔曼方程写为紧凑的矩阵形式。
定义：
\begin{itemize}
  \item 状态价值向量 $\boldsymbol{v}_\pi = \left[ v_\pi(s_1), v_\pi(s_2), \dots, v_\pi(s_n) \right]^T \in \mathbb{R}^n$；
  \item 策略即时奖励向量 $\boldsymbol{r}_\pi = \left[ r_\pi(s_1), r_\pi(s_2), \dots, r_\pi(s_n) \right]^T \in \mathbb{R}^n$，其中 $r_\pi(s) = \sum_a \pi(a|s) r(s,a)$；
  \item 策略状态转移矩阵 $\boldsymbol{P}_\pi \in \mathbb{R}^{n \times n}$，元素 $(\boldsymbol{P}_\pi)_{i,j} = \sum_a \pi(a \mid s_i) p(s_j \mid s_i,a)$。
\end{itemize}
则贝尔曼方程的矩阵形式为：
\begin{equation} \label{eq:bellman_matrix}
  \boldsymbol{v}_\pi = \boldsymbol{r}_\pi + \gamma \boldsymbol{P}_\pi \boldsymbol{v}_\pi.
\end{equation}

\subsection*{4. 策略评估：贝尔曼方程的求解}
策略评估（Policy Evaluation）即给定固定策略 $\pi$，求解贝尔曼方程以得到对应状态价值函数 $v_\pi$。

\subsubsection*{解法1：解析闭式解}
对矩阵形式直接代数整理求解：
\[
  \boldsymbol{v}_\pi = \left( \boldsymbol{I} - \gamma \boldsymbol{P}_\pi \right)^{-1} \boldsymbol{r}_\pi,
\]
其中 $\boldsymbol{I}$ 为 $n$ 阶单位矩阵。该方法仅适用于小规模状态空间，实际场景中状态数量巨大时，矩阵求逆计算复杂度极高、不可行。

\subsubsection*{解法2：迭代法（含收敛性证明）}
通过迭代更新逐步逼近真实价值函数，迭代公式为：
\[
  \boldsymbol{v}_{k+1} = \boldsymbol{r}_\pi + \gamma \boldsymbol{P}_\pi \boldsymbol{v}_k.
\]
\begin{itemize}
  \item \textbf{初始值}：可任意初始化 $\boldsymbol{v}_0$；
  \item \textbf{收敛性}：当迭代步数 $k \to \infty$ 时，$\boldsymbol{v}_k \to \boldsymbol{v}_\pi$。
\end{itemize}

\begin{proofbox}
  \textbf{迭代法收敛性证明（压缩映射原理）}
  
  设 $\boldsymbol{v}_\pi$ 为方程 \eqref{eq:bellman_matrix} 的唯一解（即真值），定义第 $k$ 步的误差向量 $\boldsymbol{\delta}_k = \boldsymbol{v}_k - \boldsymbol{v}_\pi$。

  由迭代公式与不动点方程相减，可得误差的递推关系：
  \[
    \boldsymbol{\delta}_{k+1} = \boldsymbol{v}_{k+1} - \boldsymbol{v}_\pi
    = \left(\boldsymbol{r}_\pi + \gamma \boldsymbol{P}_\pi \boldsymbol{v}_k\right)
    - \left(\boldsymbol{r}_\pi + \gamma \boldsymbol{P}_\pi \boldsymbol{v}_\pi\right)
    = \gamma \boldsymbol{P}_\pi \boldsymbol{\delta}_k .
  \]

  对上式两边同时取无穷范数（$\ell_\infty$ 范数），并结合范数的相容性（次乘法性）：
  \[
    \| \boldsymbol{\delta}_{k+1} \|_\infty
    = \| \gamma \boldsymbol{P}_\pi \boldsymbol{\delta}_k \|_\infty
    \le \gamma \cdot \| \boldsymbol{P}_\pi \|_\infty \cdot \| \boldsymbol{\delta}_k \|_\infty .
  \]

  由于 $\boldsymbol{P}_\pi$ 是行随机矩阵（每行元素非负且和为 1），其诱导的无穷范数（即最大行和）为：
  \[
    \| \boldsymbol{P}_\pi \|_\infty = \max_{i} \sum_{j} |(\boldsymbol{P}_\pi)_{i,j}| = 1 .
  \]

  结合折扣因子 $\gamma \in [0,1)$，可得误差范数的递推不等式：
  \[
    \| \boldsymbol{\delta}_{k+1} \|_\infty \le \gamma \| \boldsymbol{\delta}_k \|_\infty .
  \]

  递推展开至初始步：
  \[
    \| \boldsymbol{\delta}_k \|_\infty \le \gamma^k \| \boldsymbol{\delta}_0 \|_\infty .
  \]

  因为 $0 \le \gamma < 1$，所以当 $k \to \infty$ 时，$\gamma^k \to 0$，从而有：
  \[
    \lim_{k \to \infty} \| \boldsymbol{\delta}_k \|_\infty = 0 .
  \]

  这意味着向量序列 $\{\boldsymbol{v}_k\}$ 在范数意义下收敛于 $\boldsymbol{v}_\pi$，即：
  \[
    \lim_{k \to \infty} \boldsymbol{v}_k = \boldsymbol{v}_\pi .
  \]
  证毕。\hfill $\square$
\end{proofbox}

\begin{remarkbox}
  上述证明本质是\textbf{压缩映射（Banach 不动点）定理}在有限维向量空间的具体应用。算子 $\mathcal{T}_\pi(\boldsymbol{v}) = \boldsymbol{r}_\pi + \gamma \boldsymbol{P}_\pi \boldsymbol{v}$ 是一个 $\gamma$-收缩映射，保证了迭代的线性收敛速度，收敛速率由 $\gamma$ 决定（$\gamma$ 越接近 0，收敛越快）。
\end{remarkbox}

\subsection*{5. 动作价值函数}
动作价值函数定义为：在状态 $s$ 下执行动作 $a$，之后始终遵循策略 $\pi$ 交互所能获得的期望折扣回报：
\[
  q_\pi(s,a) = \mathbb{E}_\pi\left[ G_t \mid S_t = s, A_t = a \right].
\]
展开后的递归形式为：
\[
  q_\pi(s,a) = r(s,a) + \gamma \sum_{s' \in \mathcal{S}} p(s' \mid s,a) \, v_\pi(s').
\]

状态价值函数与动作价值函数满足如下关系：
\[
  v_\pi(s) = \sum_{a \in \mathcal{A}(s)} \pi(a \mid s) \, q_\pi(s,a).
\]

\end{document}